% !PRINTED 638
\weberpage{638}
Nach dem Satze von \gesperrt{Lagrange} (§.155) kann man jede Function
$\Theta_\nu$, der die vier Eigenschaften $\alpha)$, $\beta)$, $\gamma)$,
$\delta)$ zukommen, rational im Körper $\Omega$ der metacyklischen
Functionen von $x$ durch eine von ihnen, $\omega_\nu$, ausdrücken, wenn
nur die $\omega_0, \omega_1 \ldots \omega_{n-2}$ von einander verschieden
sind; man kann also für diese Functionen die $f_\nu$ wählen.

Denn bedeutet $u$ eine Variable, und ist

\weq{19}{(u - \omega_0)(u - \omega_1) \cdots (u - \omega_{n-2}) = \varphi(u),}

so gestattet die Function $\varphi(u)$ die Permutationen $s, t$ und auch
die Substitution $\sigma = (\varepsilon, \varepsilon^g)$; und ist also
nach dem Satze 3. des §.183 eine ganze Function von $u$ mit in Bezug auf
$x$ metacyklischen Co\"efficienten und unabhängig von $\varepsilon$.

Ist nun also $\Theta_0, \Theta_1 \ldots \Theta_{n-2}$ ein
Functionensystem, dem die Eigenschaften $\alpha) \beta) \gamma) \delta)$
zukommen, so ist die Summe

\weq{20}{\sum_{0,\, n-2}^{\nu} \frac{\Theta_\nu\, \varphi(u)}{u - \omega_\nu} = \chi(u)}

eine ganze Function $(n-2)^{\text{ten}}$ Grades von $u$, die gleichfalls
die Substitutionen $s, t, \sigma$ gestattet und die also Coëfficienten in
$\Omega$ hat. Setzt man

\weq{21}{\frac{\chi(u)}{\varphi'(u)} = \Theta(u)}

und setzt dann in (20) $u = \omega_\nu$, so folgt

\[
\Theta_\nu = \Theta(\omega_\nu),
\]

worin $\Theta$ eine rationale Function bedeutet, deren Coëfficienten
metacyklische Functionen der $x$ sind, und $\varepsilon$ nicht mehr
enthalten.

\begin{center}
§. 185.

\gesperrt{Wurzeln metacyklischer Gleichungen.}
\end{center}

Es sollen jetzt in den Formeln des vorigen Paragraphen für die Variablen
$x_0, x_1 \ldots x_{n-1}$ die Wurzeln $\xi_0, \xi_1 \ldots \xi_{n-1}$
einer in irgend einem Körper $\mathfrak{K}$
% [resolved 2026-08-06, 5x zoom strip 5/5] Fraktur K confirmed. The glyph is unmistakable at
% zoom and appears twice on this page -- "in irgend einem Koerper K irreduciblen" and
% "rational (in K) sind" -- both the same letter. \mathfrak{K} stands.
irreduciblen metacyklischen Gleichung vom Grade $n$ eingeführt werden,
und zwar in der Reihenfolge, dass die metacyklischen Functionen der
$\xi$ rational (in $\mathfrak{K}$) sind, was nach den Sätzen des vorigen
Abschnittes immer möglich ist. Wir machen aber dabei zunächst noch zwei
beschränkende Vor-
%CONT>>
